% ============================================================
%  Bagian 8 – Lampiran (Form Penggunaan Generative AI)
% ============================================================

\clearpage
\section{Lampiran}
\subsection{Form Penggunaan \textit{Generative} AI}

\begin{center}
    {\large\bfseries SURAT PERNYATAAN PENGGUNAAN AI\par}
\end{center}
\vspace{-0.3em}
\noindent Dengan ini, \sout{saya}/kami*:

\begin{tabular}{@{}l@{\hspace{1.2em}}l@{\hspace{0.8em}}p{0.62\textwidth}@{}}
    Nama/NIM           & : & Nicholas Wise Saragih Sumbayak/13524037, Nathan E C Marpaung/13524062,
                             Kurt Mikhael Purba/13524065, Aurelia Jennifer Gunawan/13524089,
                             Amanda Aurellia Salsabila/13524131 \\
    Fakultas/Sekolah   & : & STEI-K \\
    Kode Mata Kuliah   & : & IF3170 \\
    Nama Mata Kuliah   & : & Inteligensi Artifisial \\
    Judul Tugas/Proposal & : & Tugas Besar 2 \\
\end{tabular}

\vspace{0.6em}
\noindent Menyatakan bahwa \sout{saya}/kami* \underline{menggunakan/tidak menggunakan}* AI dalam pengerjaan
maupun penyusunan luaran tugas pada mata kuliah yang tertulis di atas.

\vspace{0.4em}
\noindent Jika Ya, maka alat AI/kecerdasan buatan yang digunakan adalah:

\vspace{0.4em}
\begin{table}[H]
    \centering
    \small
    \renewcommand{\arraystretch}{1.15}
    \begin{tabularx}{\textwidth}{|>{\raggedright\arraybackslash}X|c|>{\centering\arraybackslash}p{2.6cm}|}
        \hline
        \multicolumn{1}{|c|}{\textbf{Lingkup Pekerjaan}} & \textbf{Digunakan?} & \textbf{Tingkat Penggunaan} \\
        \hline
        \textbf{Pemeriksaan Ejaan:} menggunakan tools seperti Grammarly, Paperpal, DeepL,
        Quillbot, Grammarbot, LanguageTool, ProWritingAid, ChatGPT, Google Gemini,
        atau sejenisnya. & & \\
        \hline
        \textbf{Pembuatan Teks:} menggunakan tools seperti ChatGPT, Gemini, Paperpal,
        Copilot, GrammarlyGO, WordAI, WriteSonic, Jasper, Jenni AI, atau sejenisnya. & & \\
        \hline
        \textbf{Bantuan Diskusi dalam Penyusunan Konten:} menggunakan tools seperti
        ChatGPT, Gemini, Copilot, Perplexity, Jenni AI, Zoom-Companion, atau sejenisnya. & & \\
        \hline
        \textbf{Pembuatan Gambar, Video, dan/atau Grafik:} menggunakan tools seperti
        Craiyon, DALL-E, Midjourney, Stable Diffusion, Microsoft Designer, Gemini, Canva AI,
        atau sejenisnya. & & \\
        \hline
        \textbf{Lainnya, Jelaskan:}\vspace{1.6em} & & \\
        \hline
    \end{tabularx}
\end{table}

\noindent \sout{Saya}/Kami* juga menyatakan bahwa setiap penggunaan AI/kecerdasan buatan yang
dilakukan pada mata kuliah tersebut telah dinyatakan di dalam surat ini.

\vspace{0.5em}
\hfill\begin{minipage}{6.5cm}
    \raggedleft
    Bandung, \rule{3cm}{0.4pt}\par
    \vspace{1.8cm}
    \rule{5cm}{0.4pt}\par
    Nama dan NIM
\end{minipage}

